% Source-derived chapter 14: Geometric examples
% Original source: standard-conjectures-abelian-fourfolds
\chapter{Geometric examples}
\label{standard-conjectures-abelian-fourfolds-chapter-14}
\source{standard-conjectures-abelian-fourfolds}

\begin{remark}[Source section]
\label{standard-conjectures-abelian-fourfolds-chapter-14-source}
\source{standard-conjectures-abelian-fourfolds}
\source{standard-conjectures-abelian-fourfolds}
This chapter reproduces the corresponding section of the retrieved
LaTeX source, including its definitions, statements, examples, and
proofs. It has not yet been translated into Lean.
\notready
\end{remark}

% The source section title was: Geometric examples
\label{geometricexamples}
\source{standard-conjectures-abelian-fourfolds}
In this section we discuss several examples to which Theorems \ref{thmscht} and \ref{mainthm} apply non-trivially. We are  particularly interested in exotic classes on abelian fourfolds (Definition \ref{definizione}). The main result is Proposition \ref{propexemple} where we discuss the existence of exotic classes that cannot be lifted to algebraic classes in characteristic zero. The techniques of construction there are inspired by \cite{LeOo} and
 \cite{Zah}. 
 
  Other examples of exotic classes will be found  in Remark \ref{altreesotiche}. We end the section with an example (other than abelian fourfolds) for which the standard conjecture of Hodge type holds true via Theorem \ref{mainthm}.
 

\begin{proposition}
\source{standard-conjectures-abelian-fourfolds}
\notready\label{propexemple}
\source{standard-conjectures-abelian-fourfolds}
Let $p$ be a prime number and  let $K$ be an  imaginary quadratic number field where $p$ does not split. Then there exists an abelian fourfold $A$ over $\overline{\mathbb{F}}_p$ verifying the following properties.
\begin{enumerate}
\item The endomorphism  algebra is a number field that can be written as the compositum \[\End (A)_\Q = K \cdot R\] where $R$ is a totally real number field such that $[R:\Q]=4$.
(In particular $A$ is simple and has a unique CM-structure).
\item Consider the motivic decomposition from Proposition \ref{primadecomposizione}(\ref{decomposizioneQ}). Among the factors $\mathcal{M}_I$ of $\mathfrak{h}^4(A)$ there exists a (unique) factor $M$  such that $M(2)$ does not contain any Lefschetz class (Definition \ref{defexotic}) but it is Frobenius invariant (for a model of $A$ defined over a finite field).
\item\label{cond2} For any CM-lifting  of $A$ to characteristic zero (see Theorem \ref{HondaTate} and Corollary \ref{lifting motives}) the Hodge structure $R_B(M)$ is of type $(3,1),(1,3)$.
\source{standard-conjectures-abelian-fourfolds}
\item\label{cond1} For any CM-lifting  of $A$ to characteristic zero we have the equality $\End_{\NUM(\C)_\Q}  (M) = K.$ 
\source{standard-conjectures-abelian-fourfolds}
\end{enumerate}
\end{proposition}
\begin{proof}
The proof is decomposed in a series of lemmas. The final step is Lemma \ref{lemmafine}.
\end{proof}
\begin{remark}
\source{standard-conjectures-abelian-fourfolds}
\notready\label{bellaciao} Let us make some comments on the above proposition.
\source{standard-conjectures-abelian-fourfolds}
\begin{enumerate}
\item If (a model of) the fourfold $A$ verifies the Tate conjecture, then $M$ is an exotic motive in the sense of Section \ref{sectionexotic}. In particular, in each characteristic, there should be infinitely  many non-isogenous abelian fourfolds having exotic motives. 

If the Tate conjecture was highly false and no such exotic classes existed then Theorem \ref{thmscht} would follow directly from the arguments in Section \ref{sectionlefschetz}.

The Tate conjecture and the standard conjecture of Hodge type   should be thought as two independent and different  problems. The first is about the construction of algebraic classes, the second is about how they intersect (independently whether there are a lot of algebraic classes or not). It seems likely that a solution of one problem does not imply a solution for the other. For example the proof of the Tate conjecture for divisors on abelian variety  \cite[Theorem 4]{Tate} does not imply the standard conjecture of Hodge type for divisor on abelian variety (which is known by a different argument).


\item Because of the Hodge types in part (\ref{cond2}), the (expected) algebraic classes in positive characteristic cannot be lifted to algebraic classes in characteristic zero.


\item\label{unamattina} Notice that the field $F=K \otimes_\Q \Q_p$ can be any quadratic extension of $\Q_p$. This field $F$ coincides with the one in Proposition \ref{propendo}. This shows that the different cases studied in Sections \ref{p-adic periods} and \ref{computation} were needed.
\source{standard-conjectures-abelian-fourfolds}
\item The hypothesis that $K$ is totally imaginary is necessary. If $K$ were a real quadratic  number field,  conditions    (\ref{cond2}) and (\ref{cond1}) in Proposition  \ref{propexemple} would not be compatible (see the proof of  Lemma \ref{bellaidea}).
\end{enumerate}
\end{remark}

\begin{lemma}
\source{standard-conjectures-abelian-fourfolds}
\notready\label{marx}
\source{standard-conjectures-abelian-fourfolds}
Let $p$ and $K$ be as in Proposition \ref{propexemple}. There exists a totally real number field $R$  such that the following holds:
\begin{enumerate}
\item  The prime $p$ does not split in $R$.
\item The degree $[R:\Q]$ is four.
\item The field $K \otimes_\Q \Q_p$ is embeddable in the field $R \otimes_\Q \Q_p$. 
\item If $\tilde{R}$ is the normal closure of $R$ over $\Q$, then $\Gal(\tilde{R}/\Q)=\mathfrak{S}_4$. In particular we have $\Gal(\tilde{R}/R)=\mathfrak{S}_3$.
\end{enumerate}
\end{lemma}
\begin{proof}
Same as in \cite[\S 3]{LeOo}.
\end{proof}

\begin{lemma}
\source{standard-conjectures-abelian-fourfolds}
\notready\label{lenin}
\source{standard-conjectures-abelian-fourfolds}
Let $R$ be as in the above lemma. The following holds:
\begin{enumerate}
\item  The subfields of the compositum $R\cdot K$ are $\Q,K,R$ and $R\cdot K$.
\item The prime $p$ factorises in $R\cdot K$ as 
\[p=(\mathfrak{p}\cdot \bar{\mathfrak{p}})^e,\]
where $\mathfrak{p}$ and $\bar{\mathfrak{p}}$ are  two prime ideals which are exchanged by  complex conjugation and $e$ is  the ramification index.
\end{enumerate}
\end{lemma}
\begin{proof}
Let $\tilde{R}$ be as in Lemma \ref{marx}. The equality $\Gal(\tilde{R}/\Q)=\mathfrak{S}_4$ implies
\[\Gal(\tilde{R}\cdot K/\Q)=\mathfrak{S}_4\times \Z/2\Z\]
and similarly $\Gal(\tilde{R}/R)=\mathfrak{S}_3$ implies $\Gal(\tilde{R}\cdot K/R \cdot K)=\mathfrak{S}_3.$
We deduce that the subfields of $R\cdot K$ are in bijection with the subgroups of $\mathfrak{S}_4\times \Z/2\Z$ containing $\mathfrak{S}_3$. Those are precisely $\mathfrak{S}_3 , \mathfrak{S}_4, \mathfrak{S}_3\times \Z/2\Z$ and $\mathfrak{S}_4\times \Z/2\Z$, which implies (1).

Note that  the equality $R\cdot K \cong R \otimes_\Q K$ holds. Hence we have
\[(R\cdot K) \otimes_\Q \Q_p \cong  (K \otimes_\Q \Q_p) \otimes_{\Q_p} (R \otimes_\Q \Q_p) \cong (R \otimes_\Q \Q_p)^{\Gal(K \otimes_\Q \Q_p/\Q_p) }, \]
where the last equality comes from  Lemma \ref{marx}(3). As ${R \otimes_\Q \Q_p}$ is a field (by Lemma \ref{marx}(1)), the prime $p$ factorizes in $R\cdot K$ as the product of two different primes. Moreover, those two primes are exchanged by $\Gal(K \otimes_\Q \Q_p/\Q_p)$. As complex conjugation generates this Galois group it exchanges these two prime ideals. In particular they must have the same ramification index. This concludes part (2).
\end{proof}

\begin{lemma}
\source{standard-conjectures-abelian-fourfolds}
\notready\label{engels}
\source{standard-conjectures-abelian-fourfolds}
Consider the prime ideals $\mathfrak{p}$ and $\bar{\mathfrak{p}}$ of  $R\cdot K$ from the above lemma. Then there exist an integer $n$, a $p$-power $q$ and a $q$-Weil number $\alpha\in R\cdot K$ verifying the following properties:
\begin{enumerate}
 \item The ideal generated by $\alpha$ factorizes as
\[(\alpha)=\mathfrak{p}^n\cdot \bar{\mathfrak{p}}^{3n}.\]
\item For any positive integer $s$, we have the equality
\[\Q(\alpha^s)= R\cdot K.\]
\item The norm $N_{R\cdot K/K}(\alpha)\in K$ equals $q^2$.
\end{enumerate}
\end{lemma}
\begin{proof}
Consider the ideal $I=(\mathfrak{p}\cdot \bar{\mathfrak{p}}^{3})^e $. We have that $I \cdot \bar{I}= p^4$ by Lemma \ref{lenin}(2). Actually, for any $g\in\Gal(\tilde{R}\cdot K/\Q),$ we have the equality ${g(I) \cdot \overline{g(I)}= p^4}$ of ideals in $\tilde{R}\cdot K$. This follows from the case $g=\id$ together with the explicit description of the Galois group in the proof of the lemma above (which implies that complex conjugation is in the center of the Galois group). Hence, we can apply  \cite[Lemma 1]{Honda} and deduce that there exists a  $q$-Weil number $\alpha\in R\cdot K$ verifying (1).





Let us show (2). Because of Lemma \ref{lenin}(1), this amounts to showing that $\alpha^s$ does not belong to $K$ nor to $R$. Now, if  $\alpha^s$ belongs to $K$ (or to $R$) then its factorization in $R\cdot K$ would have the same exponent in $\mathfrak{p}$ and in $\bar{\mathfrak{p}}$ because there is only one prime above $p$ in $K$ (or in $R$); see the hypothesis on $K$ in Proposition \ref{propexemple} (respectively by Lemma \ref{marx}(1)).


Let us now show (3). If we compute the  norm of the ideal $(\alpha)$ we obtain
\[ (N_{R\cdot K/K} (\alpha))= N_{R\cdot K/K} (\mathfrak{p}^n\cdot \bar{\mathfrak{p}}^{3n}) = N_{R\cdot K/K} (\mathfrak{p})^n\cdot N_{R\cdot K/K} ( \bar{\mathfrak{p}})^{3n}= \tilde{p}^m,\]
where $m$ is an integer and $\tilde{p}$ is the only prime ideal above $p$ in $K$. Hence, after possibly replacing $\alpha$ by a power, we obtain that the ideal $N_{R\cdot K/K} (\alpha)$ is generated by a power of $p$. For weight reasons we have
\[(N_{R\cdot K/K} (\alpha)) = (q^2).\]
This is equivalent to the relation $N_{R\cdot K/K} (\alpha) = \xi \cdot q^2$, where $\xi$ is an invertible element of the ring of integers of $K$. As the group of  invertible elements of the ring of integers of an  imaginary quadratic field is finite, after replacing $\alpha$ by a power we get (3). Notice that such a power of $\alpha$ will still have properties (1) and (2).
\end{proof}

\begin{lemma}
\source{standard-conjectures-abelian-fourfolds}
\notready\label{elche}
\source{standard-conjectures-abelian-fourfolds}
Let $\alpha$ be a $q$-Weil number verifying the properties as in the above lemma. Let $A$ be an abelian variety over $\mathbb{F}_{q}$ whose isogeny class  corresponds to $\alpha$ under the Honda--Tate correspondance \cite{TateBour}. Then the following holds:
\begin{enumerate}
\item  The dimension of $A$ is four.
\item  $A$ is geometrically simple.
\item $\End (A)_\Q = \End (A_{\bar{\mathbb{F}}_p})_\Q = K \cdot R.$
\item The slopes of $A$ are $(1/4,3/4)$
\item There are Frobenius invariant classes in $H_\ell^4(A)$ which are not of Lefschetz type.
\end{enumerate}
\end{lemma}
\begin{proof}
By \cite{Tate} the division algebra $ \End (A)_\Q$ has center equal to $\Q(\alpha)$ which is  $R\cdot K$ by Lemma \ref{engels}. Moreover, by \cite[Theorem 1]{TateBour}, this division algebra  splits at every place except possibly at  the places $\mathfrak{p}$ and $\bar{\mathfrak{p}}$ above $p$. The local invariants there are computed by the formula in \cite[Theorem 1]{TateBour}, which gives
\[\inv_\mathfrak{p}( \End (A)_\Q)=\frac{v_\mathfrak{p}(\alpha)}{v_\mathfrak{p}(q)}\cdot [(R\cdot K)_\mathfrak{p}:\Q_p] \hspace{0.5cm} \mod \Z\]
and similarly for $\bar{\mathfrak{p}}$. 

We claim that these local invariants are trivial as well. Indeed, using the factorisation in Lemma \ref{engels}(1), we deduce that $\frac{v_\mathfrak{p}(\alpha)}{v_\mathfrak{p}(q)}=\frac{1}{4}.$ On the other hand, the degree $[(R\cdot K)_\mathfrak{p}:\Q_p]$ equals four, because  $[R\cdot K:\Q]=8$ and the two primes above $p$ are exchanged by  complex conjugation (Lemma \ref{lenin}(2)). Altogether we have that $\inv_\mathfrak{p}( \End (A)_\Q)=0$ and similarly one shows $\inv_{\bar{\mathfrak{p}}}( \End (A)_\Q)=0$ 

Because all the invariants of the $(R\cdot K)$-central algebra $ \End (A)_\Q$ are trivial, we have 
$R\cdot K= \End (A)_\Q.$  As ${[R\cdot K:\Q]=8}$  we deduce (1). 

Consider now the abelian variety $A_s$  over $\mathbb{F}_{q^s}$ whose isogeny class  corresponds to $\alpha^s$. Following the Honda--Tate correspondance $A_s$ is a simple factor of $A \times_{\mathbb{F}_q} \mathbb{F}_{q^s} $. On the other hand, all the arguments above work  by replacing $\alpha$ by  $\alpha^s$, because of  Lemma \ref{engels}(2). In particular, $A_s$  has also dimension four and $R\cdot K= \End (A_s)_\Q.$ This implies (2) and (3).

One slope has already been computed, namely $\frac{v_P(\alpha)}{v_P(q)}=\frac{1}{4}.$ Duality implies that there is also the slope $3/4$. As $A$ has dimension four there are no more slopes.

Let us  now show (5). The existence of a class such as the ones claimed is equivalent to the existence of a set $I$ consisting of four Galois conjugates of $\alpha$  whose product equals $q^2$  and such that $I$ is not stable under the action of   complex conjugation (see the proof of  Lemma \ref{musica} or \cite[\S 2]{Zah}). We claim that the relation
\[N_{R\cdot K/K}(\alpha) = q^2 \]
(Lemma \ref{engels}(3))  gives precisely the existence of those four Galois conjugates. Indeed, let $\tilde{R}$ be the normal closure of $R$ over $\Q$, by definition we have
\[N_{R\cdot K/K}(\alpha) = \prod_{g \in \Hom_K( R\cdot K , \tilde{R}\cdot K)} g(\alpha).\]
Hence it is enough to show that  complex conjugation does not stabilize the set $J=\{g(\alpha)\}_{g \in \Hom_K( R\cdot K , \tilde{R}\cdot K)}.$ As the set   $J$ is of size four and the total Galois orbit of $\alpha$ is of size $8$ there is an element of $\Gal(\tilde{R}\cdot K/\Q)$ which does not stabilize $J$.
On the other hand, thanks to the equality \[\Gal(\tilde{R}\cdot K/\Q) = \Gal(\tilde{R}\cdot K/K) \times \Gal(\tilde{R}\cdot K/\tilde{R})\]
we have that the total Galois group $\Gal(\tilde{R}\cdot K/\Q)$ is generated by its subgroup $ \Gal(\tilde{R}\cdot K/K) $ and  complex conjugation. As  $ \Gal(\tilde{R}\cdot K/K) $ stabilizes $J$,   complex conjugation cannot stabilize it.
\end{proof}
\begin{lemma}
\source{standard-conjectures-abelian-fourfolds}
\notready\label{lemmafine}
\source{standard-conjectures-abelian-fourfolds}
Let $A$ be an abelian fourfold which satisfies the properties of the lemma above. Then it also satisfies all the conditions of Proposition \ref{propexemple}.
\end{lemma}
\begin{proof}
Part (1) has already been showed. Part (2) follows from Lemma \ref{elche}(5). (Unicity comes from Lemma \ref{lemmafour}.)

Let us now show part (3). Write $\alpha, \beta, \gamma, \delta, q/\alpha, q/\beta, q/\gamma, q/ \delta $ for the eight (distinct) Frobenius eigenvalues and consider the decomposition in eigenlines
\[\mathfrak{h}^1(A)=V_\alpha \oplus V_\beta \oplus V_\gamma \oplus V_\delta \oplus V_{q/\alpha} \oplus V_{q/\beta} \oplus V_{q/\gamma} \oplus V_{q/ \delta} \]
as in Proposition \ref{decomposition}. Among these eight eigenvalues, four have slope $1/4$ and four have slope $3/4$ and they are exchanged by complex conjugation.

Fix a CM-lifting (Theorem \ref{HondaTate}). The above decomposition in eigenlines will lift as well (Corollary \ref{lifting motives}). Among the eight lines, four will belong to $H^{1,0}$ and four will belong to $H^{0,1}$ and again they are exchanged by complex conjugation. The Shimura--Taniyama  formula  \cite[Lemma 5]{TateBour}  implies that there is exactly one eigenvalue, call it $\alpha$, whose slope is $1/4$ and such that
$V_\alpha \subset H^{1,0}$. Equivalently,  there is exactly one eigenvalue, namely $q/\alpha$, whose slope is $3/4$ and such that
$V_{q/\alpha} \subset H^{0,1}$. 

Now decompose $M= M_{\alpha, \beta, \gamma, \delta} \oplus M_{q/\alpha, q/\beta, q/\gamma, q/ \delta}$ via Proposition \ref{primadecomposizione}. (After possibly renaming the eigenvalues.)
With this notation we have the relation 
\[\alpha\cdot \beta \cdot \gamma \cdot \delta=q^2.\]
By looking at the $p$-adic valuation we deduce that, among $\beta, \gamma, \delta$ there is exactly one eigenvalue of slope $1/4$, say $\beta$.  Hence  we have $V_\beta \subset H^{0,1}$ and  $V_\gamma, V_\delta \subset H^{1,0}$. Altogether we deduce
\[V_\alpha \otimes V_\beta \otimes V_\gamma \otimes V_\delta \subset H^{3,1}\]
which gives (3).

\

Let us now show part (4). By construction we can find a quadratic number field $F\subset \tilde{R}\cdot K$   such that the motive $M$  decomposes in the category $\CHM (k)_{F}$ into a sum 
\[\mathcal{M}_{I} = M_{I} \oplus  M_{\bar{I}}  \] of two motives of rank one (see Proposition \ref{primadecomposizione}). We first claim that such a field $F$ must be imaginary. If $F$ were contained in $\R$ then the Betti realization of the lifting of $M_{I}$ would respect the Hodge symmetry. As it is one dimensional for weight reasons it would be of type $(2,2)$. This contradicts part (3).

By \cite{Jann},  $D=\End_{\NUM(\C)_\Q}  (M)$  is a division algebra. By construction, $F$ splits $D$. We claim that $D=F$. Otherwise we would have $D \otimes F \cong M_{2 \times 2}(F)$ which would imply that $M_{I}$ and $M_{\bar{I}} $ are isomorphic as numerical motives. As homological and numerical equivalence is known to coincide for complex abelian varieties \cite{Lieb}, this would imply that their Betti realization are isomorphic, which is impossible because of the different Hodge types.

In conclusion,  $\End_{\NUM(\C)_\Q}  (M)$ is an imaginary quadratic field contained in $ \tilde{R}\cdot K$. On the other hand,  there is only one such field (namely $K$) because of the description of  $\Gal (\tilde{R}\cdot K/\Q)$ in the proof of Lemma \ref{lenin}.
\end{proof}
\begin{remark}
\source{standard-conjectures-abelian-fourfolds}
\notready\label{altreesotiche}
\source{standard-conjectures-abelian-fourfolds}
Let us comment on other examples of exotic motives coming from abelian fourfolds.
\begin{enumerate}
\item One can construct  an abelian fourfold $A$ over a finite field having an exotic motive whose lifting to $\C$ has Betti realization of type $(2,2)$. Such a condition means that  the CM-lifting of $A$  over $\C$   has Hodge classes which are not Lefschetz. This   situation (over $\C$) has been classified in \cite{MooZa}. So, any reduction modulo $p$ of their examples will give an abelian fourfold over a finite field of the desired type. (To avoid that the reduction modulo $p$ creates more Lefschetz classes, one can take an ordinary prime). 

As already pointed out, these examples are less interesting for the standard conjecture of Hodge type, see Remark \ref{boris}(\ref{duccio}).
\item\label{altreesotiche2} There are no exotic motives over $\overline{\mathbb{F}}_p$ (coming from abelian fourfolds) whose lifting  to $\C$ have Betti realization    of type $(4,0),(0,4)$. To show this, consider a model of the abelian fourfold over a finite field $\mathbb{F}_{q}$. Let $I$ be the set of Frobenius eigenvalues such that the corresponding eigenspaces are lifted into $H^{1,0}$ and $\bar{I}$ be the set of Frobenius eigenvalues such that the corresponding eigenspaces are lifted into $H^{0,1}$. If the cohomology group $H^{4,0}\oplus H^{0,4}$ becomes Frobenius invariant over $\mathbb{F}_q$, then $\prod_{\alpha\in I} \alpha = \prod_{\beta\in \bar{I}} \beta(=q^2).$ On the other hand, using the Shimura--Taniyama formula  \cite[Lemma 5]{TateBour}, we have that the $p$-adic valuation of  $\prod_{\alpha\in I} \alpha$ is greater than the one of $ \prod_{\beta\in \bar{I}} \beta$, except if all Frobenius eigenvalues have the same slopes. In this case the abelian variety would be isogenous to the forth power of a supersingular elliptic curve and hence all algebraic classes would be Lefschetz.
\source{standard-conjectures-abelian-fourfolds}
\item Having the results of Section 7 in mind, the last example  that needs to be discussed is that of an abelian fourfold with a four dimensional space of exotic classes. By Lemma \ref{lemmafour}, this reduces to an abelian fourfold over  $\mathbb{F}_{q^2}$  of the form $X\times E$, where $X$ is an abelian threefold and $E$ is a supersingular elliptic curve on which   Frobenius acts as $q \cdot \id$. Now the equations  (\ref{eq:4}) and  (\ref{eq:5}) imply that the existence of an exotic class on  $X\times E$ is equivalent to the existence of an exotic class on $X^2$. There are infinitely many such threefolds $X$, they have been classified in \cite{Zah}. 
\end{enumerate}
\end{remark}
\begin{proposition}
\source{standard-conjectures-abelian-fourfolds}
\notready\label{propfermat}
\source{standard-conjectures-abelian-fourfolds}
The standard conjecture of Hodge type holds true for Fermat's cubic fourfold
$X=\{x_0^3+\cdots +x_5^3=0\}\subset \mathbb{P}_k^5$ over any field $k$ (of characteristic different from $3$).
\end{proposition}
\begin{proof}
Let us first consider $X$ as variety over $\C$.  By \cite[Proposition 11]{Beauville} its Hodge structure decomposes as
\[H_B^*(X,\Q)=\Q(0) \oplus \Q(-1) \oplus \Q(-2)^{\oplus 21} \oplus V_B \oplus \Q(-3) \oplus \Q(-4)\]
where   $V_B$ is a $\Q$-Hodge structure of rank $2$ and of type $(3,1), (1,3)$.
As the Hodge conjecture is known for $X$ and its powers \cite{ShioC}, this decomposition holds true at the level of homological motives
\[M(X)=\one  \oplus \one(-1) \oplus \one(-2)^{\oplus 21} \oplus V \oplus \one(-3) \oplus \one(-4).\]
This implies that the primitive part of the motive  is of the form
\[\mathfrak{h}^{4,\prim}(X)= \one(-2)^{\oplus 20} \oplus^{\perp} V.\]
Note that the decomposition is orthogonal with respect to the cup product as the types of the Hodge structures   are different.
Finally, as the motive of $X$ is finite dimensional \cite[Lemma 5.2]{BLP}, this decomposition lifts to the level of Chow motives. (Alternatively, see Remark \ref{hypothesismainthm}(\ref{homvsrat}).)

Let us now work over $\overline{\mathbb{F}}_p$. (This is enough for our purpose, thanks to Proposition \ref{reducefinite}).  The positivity of the cup product on algebraic classes on the factor $\one(-2)^{\oplus 20}$ is clear as all these classes come from characteristic zero, see Remark \ref{remHR}. 
We are reduced to the study of algebraic classes on the two dimensional motive $V$.  As the characteristic polynomial of   Frobenius acting on $V$ is a rational polynomial of degree two, there are either zero or two rational solutions.  In the first case  the space of algebraic classes on $V$ is reduced to  zero hence  the standard conjecture of Hodge type holds trivially. In the second case the Fermat variety is supersingular and $V$ is spanned by algebraic cycles  \cite{ShioF}. In this case the standard conjecture of Hodge type holds true via Theorem \ref{mainthm}. (Note that there are infinitely many primes for which the non-trivial case occurs \cite[Theorem 2.10]{ShioF}). 
\end{proof}
\begin{remark}
\source{standard-conjectures-abelian-fourfolds}
\notready
Let us comment on applications and limits of Theorem \ref{mainthm}.
\begin{enumerate}
\item Theorem \ref{mainthm} cannot be applied to show the standard conjecture of Hodge type for abelian varieties of dimension at least five. Indeed, let $A$ be a simple abelian variety  of dimension $g$ and let $M_I\subset \mathfrak{h}^{2i}(A)$ be a factor as constructed in Proposition \ref{primadecomposizione}. By its very construction, the dimension of  $M_I$ is at least $g/i$ as $\Gal(\overline{\Q}/\Q)$ acts transitively on $\Sigma$, see Notation \ref{NotationCM}. Hence the rank of $M_I$ will never be two (except possibly in  middle degree).
\item It seems likely that using Theorem \ref{mainthm} one can show the standard conjecture of Hodge type for some special varieties as we did in Proposition \ref{propfermat} for  Fermat's  cubic fourfold.  On the other hand we do not know examples (other than abelian fourfolds) where this strategy applies  for a whole family of varieties   and we expect such examples to be rare\footnote{For example  the proof of Proposition \ref{propfermat} cannot apply to all cubic fourfolds as the smallest $\Q$-Hodge structure containing the $H^{1,3}$ of a cubic fourfold has in general dimension greater than two.}.  It is rather a miracle, based on the computations of Section 7, that for all abelian fourfolds only motives of rank two turn out to be significant.\end{enumerate}
\end{remark}
\end{appendices}


\bibliographystyle{alpha}	% (uses file "plain.bst")
%\bibliography{/Users/enrightward/Desktop/DISS_NEWFONT/masterbib_NF.bib}
\bibliography{masterbib}

%\bibliography{masterbib}
